\begin{explanationbox}
	\textbf{\textcolor{CExplanation}{一句话直觉}}

	等差数列的求和公式可以通过二次函数直观理解，其二次项系数由公差决定。

	\medskip
	\textbf{\textcolor{CExplanation}{核心拆解}}
	\begin{description}[style=nextline, leftmargin=2.8em, labelsep=0.5em, itemsep=0.45em]
		\item[\textbf{要点一（二次项来源）：}] 二次项系数$\frac{d}{2}$只与公差$d$有关。
		\item[\textbf{要点二（一次项组合）：}] 一次项系数$a_1 - \frac{d}{2}$由首项和公差共同决定。
	\end{description}

	\medskip
	\textbf{\textcolor{CExplanation}{几何本质（抛物线图象）}}

	$S_n$的图象是过原点的抛物线，但只有$n\in\mathbb{N}^*$对应离散点。

	\medskip
	\textbf{\textcolor{CExplanation}{代数意义（二次多项式）}}

	将$S_n$化为二次多项式后，可直接分析对称轴、开口方向等代数性质。

	\medskip
	\textbf{\textcolor{CExplanation}{考点价值（二次化应用）}}
	\begin{itemize}[leftmargin=*, itemsep=0.5em, topsep=0.4em]
		\item \textbf{考法一（反求首公）：} 已知$S_n$的二次表达式，比较系数即可得到首项$a_1$和公差$d$。
		\item \textbf{考法二（最值分析）：} 利用二次函数的顶点分析$S_n$的最大值或最小值（需注意$n$为正整数）。
	\end{itemize}

	\medskip
	\textbf{\textcolor{CExplanation}{顿悟点}}

	常数项为$0$是因为$S_0=0$：空集和为零。

	\medskip
	\textbf{\textcolor{CExplanation}{使用场景}}
	\begin{itemize}[leftmargin=*, itemsep=0.5em, topsep=0.4em]
		\item \textbf{场景一（快速识式）：} 题目给出$S_n$的多项式形式时，直接判断是否为等差数列前$n$项和。
		\item \textbf{场景二（性质探究）：} 利用二次函数的开口和对称轴讨论$S_n$的增减性和最值。
	\end{itemize}
\end{explanationbox}
